% Part: first-order-logic
% Chapter: introduction
% Section: first-order-logic

\documentclass[../../../include/open-logic-section]{subfiles}

\begin{document}

\olfileid{fol}{int}{fol}

% SEGMENT OLP-0140-S01
\olsection{முதல்தர தருக்கம்}

முறையிய தருக்கத்தை முதலில் அறிமுகப்படுத்திய பாடத்திலிருந்து முதல்தர தருக்கத்தை
நீங்கள் அறிந்திருக்க வாய்ப்புள்ளது.\footnote{உண்மையில், நீங்கள் அறிந்திருப்பதாகவே
ஏறத்தாழக் கருதுகிறோம்!{} அறியவில்லை எனில், \emph{forall x}
\citep{Magnus2021} போன்ற இன்னும் தொடக்கநிலையிலான பாடநூலைப் படித்து
மீள்பார்வையிடலாம்.} இதை ``அளவையடைத் தருக்கம்'' அல்லது ``பயனிலைத் தருக்கம்''
என்றும் நீங்கள் அறிந்திருக்கலாம். முதலாவதாக, முதல்தர தருக்கம் ஒரு முறையிய
மொழி. அதாவது, அதற்கு ஒரு குறிப்பிட்ட சொற்கோவை உண்டு; அதன் கோவைகள் அந்தச்
சொற்கோவையிலிருந்து உருவான குறித்தொடர்கள். ஆனால் ஒவ்வொரு குறித்தொடரும்
அனுமதிக்கப்படுவதில்லை. அனுமதிக்கப்பட்ட கோவைகளில் சொற்கள், !!{formula}s,
!!{sentence}s என வெவ்வேறு வகைகள் உள்ளன. முதல்தர தருக்கத்தின் !!{sentence}s
தான் நமது முதன்மை ஆர்வம்: அவை இயல்மொழி வாக்கியங்களுக்கான முறையிய இணையைத்
தருகின்றன; தருக்கவியலாளர் பொதுவாக ஆர்வம் கொள்ளும் கேள்விகளை அவற்றைப் பற்றிக்
கேட்கலாம். எடுத்துக்காட்டாக:

% SEGMENT OLP-0140-S02
\begin{itemize}
    \item $!B$ ஆனது~$!A$ இலிருந்து தருக்க முறையில் பின்வருகிறதா?
    \item $!A$ தருக்க முறையில் மெய்யா, பொய்யா, அல்லது
    வாய்ப்புநிலையானதா?
    \item $!A$ மற்றும் $!B$ சமானமானவையா?
\end{itemize}

% SEGMENT OLP-0140-S03
இக்கேள்விகள் முதன்மையாக முதல்தர தருக்கத்தின் !!{sentence}s இன்
``பொருள்'' பற்றியவை. எடுத்துக்காட்டாக, $!B$ ஆனது~$!A$ இலிருந்து தருக்க
முறையில் பின்வருகிறதா என்ற கேள்வியை ஒரு மெய்யியலாளர் பின்வருமாறு பகுப்பாய்வு
செய்வார்: $!A$ மெய்யாகவும்~$!B$ பொய்யாகவும் இருக்கும் நிலை ஏதேனும் உள்ளதா
($!B$ ஆனது~$!A$ இலிருந்து பின்வரவில்லை), அல்லது $!A$ ஐ மெய்யாக்கும் ஒவ்வொரு
நிலையும்~$!B$ ஐயும் மெய்யாக்குகிறதா ($!B$ ஆனது~$!A$ இலிருந்து பின்வருகிறது)?
ஆனால் ஒரு ``நிலை'' என்றால் என்ன என்பது இன்னும் நமக்குக் கூறப்படவில்லை---அதை
வழங்குவது \emph{பொருண்மையியலின்} பணி. முதல்தர தருக்கத்தின் பொருண்மையியல்,
மெய்யியலாளரின் உள்ளுணர்வுசார் ``நிலை'' என்ற கருத்துக்கும், ஒரு நிலையில்
!!a{sentence}~$!A$ \emph{மெய்யாக இருப்பது} என்றால் என்ன என்பதற்கும்
கணிதத் துல்லியமான மாதிரியை வழங்குகிறது. நாம் உருவாக்கவுள்ள அந்தக் கணிதத்
துல்லியமான மாதிரியை !!a{structure} என்கிறோம். ``இதில் மெய்'' என்பதைத்
துல்லியமாக்கும் உறவை \emph{நிறைவுறுத்தல்} உறவு என்கிறோம். ஆகவே
!!{sentence}s~$!A$ மற்றும் !!{structure}s~$\Struct{M}$ க்கு,
``$!A$ ஆனது~$\Struct{M}$ இல் நிறைவுறுத்தப்படுகிறது'' என்பதை (குறிகளில்:
$\Sat{M}{!A}$) வரையறுப்போம். இதைச் செய்தபின், ``இதிலிருந்து பின்வருகிறது''
அல்லது ``தருக்க முறையில் மெய்'' போன்ற மற்ற பொருண்மையியல் சொற்களுக்கும்
துல்லியமான வரையறைகளை வழங்கலாம். எடுத்துக்காட்டாக,
% SOURCE CORRECTION TA-INT-001: close the universal formula before the premise separator and remove the stray final bracket.
$\lforall[x][(!A(x) \lif !B(x))], \lexists[x][!A(x)] \Entails
\lexists[x][!B(x)]$ என்பது பொருந்துகிறதா என்பதை மீண்டும் கணிதத் துல்லியத்துடன்
இவ்வரையறைகள் தீர்மானிக்க உதவும். விடை நிச்சயமாக ``ஆம்.'' இயல்மொழி
வாக்கியங்களை முதல்தர தருக்கத்தில் குறியாக்க நீங்கள் ஏற்கெனவே பயிற்சி
பெற்றிருந்தால், ``எல்லா எறும்புகளும் பூச்சிகள்; எறும்புகள் உள்ளன; எனவே
பூச்சிகள் உள்ளன'' என்பதன் குறியாக்கமாக இதை அடையாளம் காண்பீர்கள். இது
வெளிப்படையாகச் செல்லுபடியான தருக்கவாதம்; ஆகவே நமது முறையிய மொழிக்கான
``இதிலிருந்து பின்வருகிறது'' என்பதன் கணித மாதிரியும் அதே விடையைத் தர வேண்டும்.

% SEGMENT OLP-0140-S04
முறையிய தருக்கத்தின் முதல் அறிமுகத்திலிருந்து உங்களுக்கு நினைவிருக்கக்கூடிய
மற்றொரு தலைப்பு \emph{!!{derivation}s} இருப்பது. முதல் முறையிய தருக்கப்
பாடத்தைப் படித்திருந்தால், அத்தகைய !!{derivation}s ஐக் கண்டறிய உங்கள்
ஆசிரியர் பயிற்சி அளித்திருப்பார்; மேலுள்ள உட்கிடை பொருந்துகிறது என்பதைக்
காட்டும் !!a{derivation} ஐக்கூடக் கண்டிருக்கலாம். !!{derivation}s ஐ வழங்கப்
பல வழிகள் உள்ளன: ``இயல்பான வருவித்தல்'' அல்லது ``மெய்மை மரங்கள்'' எனப்படும்
ஒன்றை நீங்கள் செய்திருக்கலாம்; இன்னும் பல முறைகளும் உள்ளன. மேலுள்ள
தருக்கவியலாளரின் கேள்விகளுக்குப் பதிலளிக்க உதவும் கருவிகளை வழங்குவதே
!!{derivation} முறைமைகளின் நோக்கம். எடுத்துக்காட்டாக,
% SOURCE CORRECTION TA-INT-002: balance the quantified premise and conclusion formulas in the derivation example.
$\lforall[x][(!A(x) \lif !B(x))]$ மற்றும் $\lexists[x][!A(x)]$ ஆகியவற்றை
முன்கூற்றுகளாகவும், $\lexists[x][!B(x)]$ ஐ முடிவாகவும் (கடைசி வரியாகவும்)
கொண்ட இயல்பான வருவித்தல் !!{derivation} ஒன்று, $\lforall[x][(!A(x) \lif
!B(x))]$ மற்றும் $\lexists[x][!A(x)]$ ஆகியவற்றிலிருந்து
$\lexists[x][!B(x)]$ தருக்க முறையில் பின்வருகிறது என்பதை \emph{சரிபார்க்கிறது}.

% SEGMENT OLP-0140-S05
ஆனால் அது ஏன்? முதல் பார்வையில், !!{derivation} முறைமைகளுக்கும்
பொருண்மையியலுக்கும் தொடர்பே இல்லை: ஒரு முறையிய !!{derivation} ஐ வழங்குவது,
சில விதிகளால் ஆளப்படும் முறைகளில் குறிகளை அடுக்குவதை மட்டுமே கொண்டது;
``நிலைகள்'' அல்லது ``இதில் மெய்'' என்பதை அவை குறிப்பிடவே இல்லை.
!!{derivation} முறைமைகளுக்கும் பொருண்மையியலுக்கும் இடையிலான தொடர்பு ஒரு
மேல்தருக்கவியல் ஆய்வால் நிறுவப்பட வேண்டும். எடுத்துக்காட்டாக,
% SOURCE CORRECTION TA-INT-003: balance the quantified premise and conclusion formulas in the metalogical biconditional.
$\lforall[x][(!A(x) \lif !B(x))]$ மற்றும் $\lexists[x][!A(x)]$ என்ற
முன்கூற்றுகளிலிருந்து $\lexists[x][!B(x)]$ ஐ வருவிக்கும் ஒரு முறையிய
!!{derivation} இருப்பது அப்போதும் அப்போது மட்டுமே
$\lforall[x][(!A(x) \lif !B(x))]$ மற்றும் $\lexists[x][!A(x)]$ ஆகியவை
சேர்ந்து~$\lexists[x][!B(x)]$ ஐ உட்கிடையச் செய்கின்றன என்ற கணித நிறுவல்
தேவை. ஆனால் இதைச் செய்வதற்கு முன், தொடரமைப்பு மற்றும் பொருண்மையியலின்
வரையறைகளைச் சரியாக அமைக்கும் பெரும் நுணுக்கமான பணியை முடிக்க வேண்டும்.

\end{document}
